% ===========================================================================
% exp05_dbscan.tex - complete code listings for Experiment 5
% Source: notebooks/02_density_based_learning.ipynb (executed)
% ===========================================================================

\begin{lstlisting}[style=mlcode,caption={Core numerical and plotting libraries}]
# ---- Core numerical and plotting libraries ----
import numpy as np
import pandas as pd
import matplotlib.pyplot as plt
import seaborn as sns

# ---- Density-based clustering algorithms ----
from sklearn.cluster import DBSCAN, HDBSCAN
# ---- Nearest-neighbor search (used to pick epsilon) ----
from sklearn.neighbors import NearestNeighbors
# ---- Evaluation metrics ----
from sklearn.metrics import silhouette_score

# Aesthetics setup (consistent look across all notebooks)
plt.style.use('seaborn-v0_8-whitegrid' if 'seaborn-v0_8-whitegrid' in plt.style.available else 'default')
plt.rcParams['font.sans-serif'] = 'DejaVu Sans'
plt.rcParams['figure.figsize'] = (11, 6)
plt.rcParams['figure.dpi'] = 110

np.random.seed(42)
print("Density-based libraries loaded successfully!")
\end{lstlisting}

\begin{lstlisting}[style=mlcode,caption={Load the real USGS earthquake catalogue (global, 2023)}]
# ---- Load the real USGS earthquake catalogue (global, 2023) ----
df = pd.read_csv('../data/earthquakes.csv')
print(f"Full catalogue (magnitude >= 4.0): {len(df):,} events")
print(df.head(3).to_string(index=False))

# Keep well-recorded moderate+ events: magnitude >= 4.5
quakes = df[df['mag'] >= 4.5].reset_index(drop=True)
print(f"\nAnalysed subset (magnitude >= 4.5): {len(quakes):,} events")

# ---- Convert epicenters to radians [latitude, longitude] ----
# scikit-learn's 'haversine' metric expects radian coordinates and computes
# great-circle distance on the sphere (correct for global data).
X_rad = np.radians(quakes[['latitude', 'longitude']].values)
print("Coordinate range (deg): lat", quakes.latitude.min(), "to", quakes.latitude.max(),
      "| lon", quakes.longitude.min(), "to", quakes.longitude.max())
\end{lstlisting}

\begin{lstlisting}[style=mlcode,caption={Choosing epsilon from the k-distance graph (haversine metric)}]
# ---- Choosing epsilon from the k-distance graph (haversine metric) ----
# Sorted distance to each event's k-th nearest neighbor. The knee marks the
# transition from dense seismic zones to sparse, isolated events.
min_samples = 10
nbrs = NearestNeighbors(n_neighbors=min_samples, metric='haversine').fit(X_rad)
distances, _ = nbrs.kneighbors(X_rad)
k_distances = np.sort(distances[:, -1])   # radians (0.03 rad = 191 km on Earth)

plt.figure(figsize=(9, 5))
plt.plot(k_distances, color='crimson', linewidth=2.5)
plt.title(f"k-Distance Graph (k={min_samples}) for Epsilon Selection", fontsize=14, fontweight='bold')
plt.xlabel("Earthquakes Sorted by Distance", fontsize=12)
plt.ylabel(f"{min_samples}-NN Distance (radians)", fontsize=12)
plt.axhline(y=0.03, color='navy', linestyle='--', label='Selected ε = 0.03 rad (~191 km)')
plt.legend()
plt.tight_layout()
plt.show()
\end{lstlisting}

\begin{lstlisting}[style=mlcode,caption={Fit DBSCAN with the chosen epsilon and MinPts}]
# ---- Fit DBSCAN with the chosen epsilon and MinPts ----
# DBSCAN grows clusters from core events, connects them via density-reachability,
# and labels everything else as noise (-1).
eps_selected = 0.03   # radians -> 0.03 * 6371 km ~ 191 km
dbscan = DBSCAN(eps=eps_selected, min_samples=min_samples, metric='haversine')
db_labels = dbscan.fit_predict(X_rad)

# Identify which events are core points (returned by the model) vs. border points
core_samples_mask = np.zeros_like(db_labels, dtype=bool)
core_samples_mask[dbscan.core_sample_indices_] = True

n_clusters_db = len(set(db_labels)) - (1 if -1 in db_labels else 0)
n_noise_db = np.sum(db_labels == -1)
print(f"DBSCAN: {n_clusters_db} clusters, {n_noise_db:,} noise events ({n_noise_db / len(X_rad) * 100:.1f}%)")

# ---- Visualize epicenters: clustered events in color, noise as black crosses ----
plt.figure(figsize=(12, 5.5))
palette = sns.color_palette("tab20", max(n_clusters_db, 1))
for k in range(n_clusters_db):
    mask = db_labels == k
    plt.scatter(quakes.longitude[mask], quakes.latitude[mask], s=6, color=palette[k % 20], alpha=0.75)

noise_mask = db_labels == -1
plt.scatter(quakes.longitude[noise_mask], quakes.latitude[noise_mask], s=6, color='black', marker='x', label='Noise (-1)')
plt.title(f"DBSCAN on Global Earthquakes 2023 (ε={eps_selected} rad, {n_clusters_db} clusters)", fontsize=13, fontweight='bold')
plt.xlabel("Longitude", fontsize=12)
plt.ylabel("Latitude", fontsize=12)
plt.legend(loc='lower left')
plt.tight_layout()
plt.show()
\end{lstlisting}

\begin{lstlisting}[style=mlcode,caption={Parameter sensitivity: DBSCAN under several (eps, min\_samples) settings}]
# ---- Parameter sensitivity: DBSCAN under several (eps, min_samples) settings ----
# The manual requires at least two settings so the effect of eps can be discussed.
print("eps (rad)   approx km   clusters   noise %")
for eps_test in [0.02, 0.03, 0.05]:
    lbl = DBSCAN(eps=eps_test, min_samples=10, metric='haversine').fit_predict(X_rad)
    n_clusters_test = len(set(lbl)) - (1 if -1 in lbl else 0)
    noise_pct = (lbl == -1).mean() * 100
    print(f"  {eps_test:.2f}       {eps_test * 6371:5.0f}      {n_clusters_test:3d}      {noise_pct:5.1f}%")
\end{lstlisting}
